\chapter{相空间计算公式}
\label{app:math}

本附录集中列出正文反复使用的变换、对应规则与排序公式。为避免量纲混淆，前两节用 $(x,p_x)$ 表示满足 $[\hat x,\hat p_x]=\ii\hbar$ 的有量纲坐标--动量；从“复变量 Bopp 对应”开始，玻色模式使用无量纲正交分量 $(Q,P)$，满足
$\alpha=(Q+\ii P)/\sqrt2$ 与 $[\hat Q,\hat P]=\ii$。两套变量的测度和 $\hbar$ 因子不能交叉套用。

\section{Weyl 符号与反变换}

对一维正则坐标，算符 $\Ohat$ 的 Weyl 符号为
\begin{equation}
  O_W(x,p_x)=\int_{-\infty}^{\infty}\dd y\,
  \ee^{-\ii p_x y/\hbar}
  \mel{x+y/2}{\Ohat}{x-y/2}.
\end{equation}
反变换写成
\begin{equation}
  \mel{x}{\Ohat}{x'}=
  \int\frac{\dd p_x}{2\pi\hbar}\,
  \ee^{\ii p_x(x-x')/\hbar}
  O_W\!\left(\frac{x+x'}2,p_x\right).
\end{equation}
密度算符的 Wigner 函数采用
\begin{equation}
  W(x,p_x)=\frac{1}{2\pi\hbar}
  \int\dd y\,\ee^{-\ii p_x y/\hbar}
  \mel{x+y/2}{\rhohat}{x-y/2},
\end{equation}
于是 $\int\dd x\dd p_x\,W=1$，且
\begin{equation}
  \qavg{\Ohat}=\int\dd x\dd p_x\,W(x,p_x)O_W(x,p_x).
\end{equation}
不同文献会把 $2\pi\hbar$ 放在 $W$ 或相空间测度中；只要正反变换、归一化和求迹公式成套一致即可。

\section{星积与 Moyal 括号}

一维星积是
\begin{equation}
  A_W\star B_W=A_W
  \exp\!\left[
  \frac{\ii\hbar}{2}
  \left(
  \overleftarrow\partial_x\overrightarrow\partial_{p_x}-
  \overleftarrow\partial_{p_x}\overrightarrow\partial_x
  \right)\right]B_W.
\end{equation}
算符乘积、对易子和反对易子分别映射为
\begin{align}
  (\hat A\hat B)_W&=A_W\star B_W,\\
  \frac{1}{\ii\hbar}[\hat A,\hat B]_W
  &=\frac{2}{\hbar}A_W\sin\!\left(\frac{\hbar\Lambda}{2}\right)B_W,\\
  \frac12\{\hat A,\hat B\}_W
  &=A_W\cos\!\left(\frac{\hbar\Lambda}{2}\right)B_W,
\end{align}
其中
$\Lambda=\overleftarrow\partial_x\overrightarrow\partial_{p_x}-
\overleftarrow\partial_{p_x}\overrightarrow\partial_x$。Moyal 括号的低阶展开为
\begin{equation}
  \{A,B\}_{\rm M}=\{A,B\}_{\rm P}
  -\frac{\hbar^2}{24}A\Lambda^3B+\order(\hbar^4).
\end{equation}

\section{复变量 Bopp 对应}

单模密度算符左右乘法对应
\begin{align}
  \ha\rhohat &\longleftrightarrow
  \left(\alpha+\frac12\partial_{\alpha^*}\right)W,&
  \had\rhohat &\longleftrightarrow
  \left(\alpha^*-\frac12\partial_\alpha\right)W,\\
  \rhohat\ha &\longleftrightarrow
  \left(\alpha-\frac12\partial_{\alpha^*}\right)W,&
  \rhohat\had &\longleftrightarrow
  \left(\alpha^*+\frac12\partial_\alpha\right)W.
\end{align}
这些规则可直接生成主方程的 Wigner 偏微分方程。使用时应把导数作用到其右侧的完整乘积，不能只作用到紧邻因子。

\section{常用单模 Weyl 符号}

\begin{center}
\begin{tabular}{ll}
\toprule
算符 & Weyl 符号 \\
\midrule
$\ha$，$\had$ & $\alpha$，$\alpha^*$ \\
$\hat n=\had\ha$ & $|\alpha|^2-1/2$ \\
$\hat n^2$ & $|\alpha|^4-|\alpha|^2$ \\
$\had{}^2\ha^2=\hat n(\hat n-1)$
& $|\alpha|^4-2|\alpha|^2+1/2$ \\
$\hat Q^2$，$\hat P^2$ & $Q^2$，$P^2$ \\
$\tfrac12(\hat Q\hat P+\hat P\hat Q)$ & $QP$ \\
\bottomrule
\end{tabular}
\end{center}
由此得到
\begin{align}
  \qavg{\hat n}&=\ensavg{|\alpha|^2}-\frac12,\\
  \qavg{\hat n(\hat n-1)}
  &=\ensavg{|\alpha|^4-2|\alpha|^2+1/2}.
\end{align}

\section{多模正规排序四次项}

对正则模式指标 $i,j,k,l$，
\begin{align}
  (\had_i\had_j\ha_k\ha_l)_W
  ={}&\alpha_i^*\alpha_j^*\alpha_k\alpha_l\\
  &-\frac12\left(
  \delta_{ik}\alpha_j^*\alpha_l+
  \delta_{il}\alpha_j^*\alpha_k+
  \delta_{jk}\alpha_i^*\alpha_l+
  \delta_{jl}\alpha_i^*\alpha_k\right)\\
  &+\frac14(\delta_{ik}\delta_{jl}+\delta_{il}\delta_{jk}).
\end{align}
令所有指标相同即得到单模四次公式；令 $i=l$、$j=k$ 且 $i\neq j$，得到
$(\hat n_i\hat n_j)_W=(|\alpha_i|^2-1/2)(|\alpha_j|^2-1/2)$。

\section{投影场公式}

有限模式场及投影核为
\begin{equation}
  \hpsi_{\mathcal C}(x)=\sum_{n\in\mathcal C}\phi_n(x)\ha_n,
  \qquad
  \delta_{\mathcal C}(x,x')=\sum_{n\in\mathcal C}
  \phi_n(x)\phi_n^*(x').
\end{equation}
常用估计器汇总为
\begin{align}
  \qavg{\hpsid(x)\hpsi(x')}
  &=\ensavg{\psi^*(x)\psi(x')}
  -\frac12\delta_{\mathcal C}(x,x'),\\
  \qavg{\hpsid(x)^2\hpsi(x)^2}
  &=\ensavg{|\psi|^4-2\delta_{\mathcal C}|\psi|^2
  +\frac12\delta_{\mathcal C}^2}.
\end{align}
接触相互作用的单分量 TWA 漂移为
\begin{equation}
  \ii\hbar\partial_t\psi_{\mathcal C}
  =\mathcal P_{\mathcal C}\left\{
  [h_0+g(|\psi_{\mathcal C}|^2-\delta_{\mathcal C}(x,x))]
  \psi_{\mathcal C}\right\}.
\end{equation}

\section{Fourier 约定速查}

对 $M$ 点周期格点，$x_j=j\Delta x$、$k_\ell=2\pi\ell/L$，使用幺正变换
\begin{equation}
  \alpha_\ell^{(k)}=\frac1{\sqrt M}\sum_j
  \ee^{-\ii k_\ell x_j}\alpha_j,
  \qquad
  \alpha_j=\frac1{\sqrt M}\sum_\ell
  \ee^{\ii k_\ell x_j}\alpha_\ell^{(k)}.
\end{equation}
于是
\begin{equation}
  \sum_j|\alpha_j|^2=\sum_\ell|\alpha_\ell^{(k)}|^2,
  \qquad
  N=\sum_\ell\left(\ensavg{|\alpha_\ell^{(k)}|^2}-\frac12\right).
\end{equation}
若程序库采用未归一化正变换和 $1/M$ 反变换，应在代码封装层显式补上 $1/\sqrt M$，而不是在每个观测量中分别猜测因子。
